% Evita que Quarto/Pandoc genere una portada automática.
% La portada institucional se construye manualmente en index.qmd.
\makeatletter
\renewcommand{\maketitle}{}
\makeatother

\usepackage{times}
\usepackage{graphicx}
\usepackage{array}
\usepackage{framed}
\usepackage{xcolor}
\usepackage{amsmath,amsfonts,amssymb}
\usepackage{float}
\usepackage{placeins}
\usepackage{subcaption}
\usepackage{algorithm}
\usepackage{algpseudocode}
\usepackage{tikz}
\usetikzlibrary{shapes.geometric, arrows, decorations.pathmorphing, shadows, fadings, circuits.ee.IEC}
\usepackage{pdflscape}
\usepackage{longtable}
\usepackage{multirow}
\usepackage{adjustbox}
\usepackage{booktabs}
\usepackage{setspace}
\usepackage{fancyhdr}
\usepackage{titlesec}
\usepackage{lmodern}
\usepackage{microtype}
\usepackage{caption}
\usepackage{nomencl}

% Nombres institucionales en español
\addto\captionsspanish{%
  \renewcommand\contentsname{Índice general}
  \renewcommand\bibname{Referencias}
  \renewcommand\tablename{Tabla nro.}
}
\captionsetup[figure]{name=Figura}

% Profundidad de numeración
\setcounter{secnumdepth}{4}

% Espaciado
\onehalfspacing
\sloppy

% Encabezados y pies de página
\pagestyle{fancy}
\fancyhf{}
\fancyhead[R]{\color{gray}{\thepage}}
\fancyhead[L]{\color{gray}{\nouppercase{\leftmark}}}
\renewcommand{\headrulewidth}{0.4pt}

% Formato de capítulos inspirado en la plantilla Overleaf original
\titleformat{\chapter}[display]
  {\normalfont\bfseries}
  {\filleft
    {\normalfont\Large\textsc{\chaptertitlename}}\hspace{10pt}
    {\centering\fontsize{20}{90}\selectfont\thechapter}}
  {10pt}
  {\titlerule[2.5pt]\vskip3pt\LARGE}

% Nomenclatura opcional
\renewcommand{\nomname}{Nomenclatura y Notaciones}
\makenomenclature

% Columnas centradas
\newcolumntype{C}[1]{>{\centering\arraybackslash}m{#1}}
